\subsection{معرفی و انگیزه}
این مکانیزم بر مبنای مبانی نظری پیشرفته در ادبیات توجه و فشرده‌سازی \cite{ref_tomoecompressor_2024} توسعه یافته است.

تبدیل لایه‌های متراکم خطی به ساختار آمیخته‌ای از خبره‌ها (\lr{Mixture of Experts - MoE}) یکی از نوآورانه‌ترین راهکارها برای حفظ ظرفیت مدل در عین افزایش کارایی محاسباتی و کاهش بار محاسباتی فعال در هر گام استنباطی است. مدل‌های سنتی ترنسفورمر حجم عظیمی از محاسبات را در لایه‌های پیش‌خوران (\lr{FFN}) و نگاشت‌های پروجکشن توجه صرف می‌کنند که در آن‌ها تمامی نورون‌ها برای هر توکن ورودی فعال می‌شوند.

موتور فشرده‌سازی \lr{ToMoE Compressor (Top-k Mixture of Experts by Weight Clustering)} با هدف بازآرایی ساختاری شبکه‌های عصبی چگال به معماری خلوت و آگاه از مسیر (\lr{Sparse Routed Architecture}) توسعه یافته است. در این الگوریتم، به جای آموزش یک مدل \lr{MoE} از نقطه صفر، وزن‌های لایه‌های خطی و \lr{Conv1D} از طریق خوشه‌بندی نورون‌های خروجی (\lr{Weight Clustering}) به چند زیرمجموعه متمایز از خبره‌ها افراز شده و یک شبکه مسیریاب پویا (\lr{Router}) وظیفه تخصیص توکن‌ها به برترین خبره‌ها (\lr{Top-k Routing}) را بر عهده می‌گیرد.
